\section{Scriptural Date and Location}
\zlabel{triptych:concise:chronology:start}
% Generated. Do not edit.
% Deterministic projection of research/chronology.toml.
% schema: 4
% document: liturgy/roman-rite/1962/propers/temporal/60-twentieth-after-pentecost
% generated-by: tools/tpt proper-chronology annotations
\newcommand{\chronologyannotationclaim}[6]{#6}
\newcommand{\chronologyannotationcomparisonclaim}[8]{#8}
\newcommand{\chronologyannotationanchorclaim}[9]{#9}
\newcommand{\chronologyannotationanchorgroup}[8]{#8}
\newcommand{\chronologyannotationreach}[3]{}
\newcommand{\chronologyannotationgroup}[3]{#3}
\newcommand{\chronologyannotationcomparisongroup}[4]{#4}
\newcommand{\chronologyannotation}[1]{%
  \ifcsname triptychchronologyannotation@#1\endcsname
    \csname triptychchronologyannotation@#1\endcsname
  \else
    \PackageError{triptych}{No chronology annotation for #1}{}%
  \fi}
% comparison-dependency: src/gpt/liturgy/roman-rite/1962/propers/temporal/60-twentieth-after-pentecost/research/chronology-profile-comparisons.toml sha256=11a4bde1a39800640daad7a03563050ab63f20b33c038c8773a288bf407b2b9d
% introit: Introit
\expandafter\def\csname triptychchronologyannotation@introit\endcsname{%
  \chronologyannotationgroup{historical-setting}{nonuniform}{\textbf{Historical setting} -- No single occasion date applies across every cited locus.}%
  \space \chronologyannotationgroup{composition}{nonuniform}{\textbf{Composition} -- No single date applies across every cited locus.}%
}
% epistle: Epistle
\expandafter\def\csname triptychchronologyannotation@epistle\endcsname{%
  \chronologyannotationgroup{composition}{disputed}{\textbf{Composition} -- disputed: \chronologyannotationreach{Eph.5.15}{Eph}{true}\chronologyannotationreach{Eph.5.16}{Eph}{true}\chronologyannotationreach{Eph.5.17}{Eph}{true}\chronologyannotationreach{Eph.5.18}{Eph}{true}\chronologyannotationreach{Eph.5.19}{Eph}{true}\chronologyannotationreach{Eph.5.20}{Eph}{true}\chronologyannotationreach{Eph.5.21}{Eph}{true}\chronologyannotationclaim{composition.epistle-to-the-ephesians}{composition}{catholic-traditional-v1}{disputed}{a period between 58 and 63}{A.D. 58--63}; \chronologyannotationreach{Eph.5.15}{Eph}{true}\chronologyannotationreach{Eph.5.16}{Eph}{true}\chronologyannotationreach{Eph.5.17}{Eph}{true}\chronologyannotationreach{Eph.5.18}{Eph}{true}\chronologyannotationreach{Eph.5.19}{Eph}{true}\chronologyannotationreach{Eph.5.20}{Eph}{true}\chronologyannotationreach{Eph.5.21}{Eph}{true}\chronologyannotationclaim{composition.epistle-to-the-ephesians}{composition}{catholic-traditional-v1}{disputed}{(Philemon; Colossians; Ephesians; Philippians), 61}{A.D. 61}.}%
  \space \chronologyannotationcomparisongroup{catholic-critical-v1}{composition}{composition-only}{\textbf{Composition (Catholic critical chronology, for comparison)} -- disputed: \chronologyannotationreach{Eph.5.15}{Eph}{true}\chronologyannotationreach{Eph.5.16}{Eph}{true}\chronologyannotationreach{Eph.5.17}{Eph}{true}\chronologyannotationreach{Eph.5.18}{Eph}{true}\chronologyannotationreach{Eph.5.19}{Eph}{true}\chronologyannotationreach{Eph.5.20}{Eph}{true}\chronologyannotationreach{Eph.5.21}{Eph}{true}\chronologyannotationcomparisonclaim{critical.ephesians-later-disciple}{composition}{catholic-critical-v1}{catholic-critical-v1}{disputed}{passage.united-states-conference-of-catholic-bishops.new-american-bible-revised-edition.english-usccb-web-2026-09-28.ephesians-introduction}{around A.D. 80--100}{Ephesians under the later-disciple hypothesis (approximate), c. A.D. 80--100}.}%
}
% gradual: Gradual
\expandafter\def\csname triptychchronologyannotation@gradual\endcsname{%
  \chronologyannotationgroup{traditional-attribution}{disputed}{\textbf{Traditional attribution} -- disputed: \chronologyannotationreach{Ps.144.15}{Ps.107 Ps.144}{true}\chronologyannotationreach{Ps.144.16}{Ps.107 Ps.144}{true}\chronologyannotationclaim{israel.monarchy.david-traditional-era}{traditional-attribution}{catholic-traditional-v1}{disputed}{reigned from 1055 to 1015 B.C.}{David (reign in the usual chronology), B.C. 1055--1015}.}%
  \space \chronologyannotationgroup{historical-setting}{research-pending}{\textbf{Historical setting} -- Occasion date unresolved.}%
  \space \chronologyannotationgroup{composition}{preferred}{\textbf{Composition}: \chronologyannotationreach{Ps.144.15}{Ps}{true}\chronologyannotationreach{Ps.144.16}{Ps}{true}\chronologyannotationclaim{critical.psalms.latest-composition-boundary}{composition}{catholic-critical-v1}{preferred}{before the Maccabean period, around 165 B.C.; no individual psalm can be dated securely}{Before c. 165 B.C.}}%
}
% alleluia: Alleluia
\expandafter\def\csname triptychchronologyannotation@alleluia\endcsname{%
  \chronologyannotationgroup{traditional-attribution}{disputed}{\textbf{Traditional attribution} -- disputed: \chronologyannotationreach{Ps.107.2}{Ps.107 Ps.144}{true}\chronologyannotationclaim{israel.monarchy.david-traditional-era}{traditional-attribution}{catholic-traditional-v1}{disputed}{reigned from 1055 to 1015 B.C.}{David (reign in the usual chronology), B.C. 1055--1015}.}%
  \space \chronologyannotationgroup{historical-setting}{research-pending}{\textbf{Historical setting} -- Occasion date unresolved.}%
  \space \chronologyannotationgroup{composition}{preferred}{\textbf{Composition}: \chronologyannotationreach{Ps.107.2}{Ps}{true}\chronologyannotationclaim{critical.psalms.latest-composition-boundary}{composition}{catholic-critical-v1}{preferred}{before the Maccabean period, around 165 B.C.; no individual psalm can be dated securely}{Before c. 165 B.C.}}%
}
% gospel: Gospel
\expandafter\def\csname triptychchronologyannotation@gospel\endcsname{%
  \chronologyannotationgroup{narrated-event}{preferred}{\textbf{Event}: \chronologyannotationreach{John.4.46}{John.4.46-53}{false}\chronologyannotationreach{John.4.47}{John.4.46-53}{false}\chronologyannotationreach{John.4.48}{John.4.46-53}{false}\chronologyannotationreach{John.4.49}{John.4.46-53}{false}\chronologyannotationreach{John.4.50}{John.4.46-53}{false}\chronologyannotationreach{John.4.51}{John.4.46-53}{false}\chronologyannotationreach{John.4.52}{John.4.46-53}{false}\chronologyannotationreach{John.4.53}{John.4.46-53}{false}\chronologyannotationclaim{life-of-christ.healing-of-the-rulers-son-at-cana}{narrated-event}{catholic-traditional-v1}{preferred}{Passover, A.U.C. 779 - about Pentecost, 780}{The healing of the ruler's son at Cana, A.D. 26--27 (derived)}.}%
  \space \chronologyannotationgroup{composition}{disputed}{\textbf{Composition} -- disputed: \chronologyannotationreach{John.4.46}{John}{true}\chronologyannotationreach{John.4.47}{John}{true}\chronologyannotationreach{John.4.48}{John}{true}\chronologyannotationreach{John.4.49}{John}{true}\chronologyannotationreach{John.4.50}{John}{true}\chronologyannotationreach{John.4.51}{John}{true}\chronologyannotationreach{John.4.52}{John}{true}\chronologyannotationreach{John.4.53}{John}{true}\chronologyannotationclaim{composition.gospel-of-john}{composition}{catholic-traditional-v1}{disputed}{from the year 90 to 100 (approximately)}{c. A.D. 90--100}; \chronologyannotationreach{John.4.46}{John}{true}\chronologyannotationreach{John.4.47}{John}{true}\chronologyannotationreach{John.4.48}{John}{true}\chronologyannotationreach{John.4.49}{John}{true}\chronologyannotationreach{John.4.50}{John}{true}\chronologyannotationreach{John.4.51}{John}{true}\chronologyannotationreach{John.4.52}{John}{true}\chronologyannotationreach{John.4.53}{John}{true}\chronologyannotationclaim{composition.gospel-of-john}{composition}{catholic-traditional-v1}{disputed}{96 or one of the succeeding years}{A.D. 96--100}.}%
}
% offertory: Offertory
\expandafter\def\csname triptychchronologyannotation@offertory\endcsname{%
  \chronologyannotationgroup{historical-setting}{preferred}{\textbf{Historical setting}: \chronologyannotationreach{Ps.136.1}{Ps.136}{true}\chronologyannotationclaim{israel.exile.psalm-136-captive-lament}{historical-setting}{catholic-traditional-v1}{preferred}{For the time of Jeremias, and the captivity of Babylon}{The Babylonian captives' lament, after Jerusalem's destruction under Sedecias (derived)}.}%
  \space \chronologyannotationgroup{historical-setting}{preferred}{\textbf{Historical setting}: \chronologyannotationreach{Ps.136.1}{Ps.136}{true}\chronologyannotationclaim{israel.exile.seventy-years}{historical-setting}{catholic-traditional-v1}{preferred}{seventy years}{The seventy years of servitude to Babylon, ``seventy years'' (duration)}.}%
  \space \chronologyannotationgroup{composition}{preferred}{\textbf{Composition}: \chronologyannotationreach{Ps.136.1}{Ps}{true}\chronologyannotationclaim{critical.psalms.latest-composition-boundary}{composition}{catholic-critical-v1}{preferred}{before the Maccabean period, around 165 B.C.; no individual psalm can be dated securely}{Before c. 165 B.C.}}%
  \space \chronologyannotationanchorgroup{israel.exile.psalm-136-captive-lament}{historical-setting}{catholic-traditional-v1}{For the time of Jeremias, and the captivity of Babylon}{israel.exile.third-captivity}{after}{context}{\textbf{Historical background: The third captivity of Juda: the destruction of Jerusalem under Sedecias}: Preferred: \chronologyannotationanchorclaim{israel.exile.third-captivity}{0}{catholic-traditional-v1}{preferred}{passage.george-leo-haydock.douay-rheims-with-haydock-commentary.2014-loreto-feeney-memorial.psalm-70-captivities-usher-chronology}{reported-excluded}{A.M. 3416}{Haydock reporting Ussher; A.M. is not converted}{A.M. 3416 (Haydock reporting Ussher; A.M. is not converted)}; alternatives: \chronologyannotationanchorclaim{israel.exile.third-captivity}{1}{catholic-traditional-v1}{alternate}{artifact.catholic-encyclopedia.volume-3.new-york-1908.newadvent-03315a-9cb437e2,artifact.catholic-encyclopedia.volume-14.new-york-1912.newadvent-14499a-82b12b6e}{traditional-catholic}{586 B.C.}{Reid and Meistermann, Catholic Encyclopedia}{B.C. 586 (Reid and Meistermann, Catholic Encyclopedia)}, \chronologyannotationanchorclaim{israel.exile.third-captivity}{2}{catholic-traditional-v1}{alternate}{artifact.catholic-encyclopedia.volume-8.new-york-1910.newadvent-08652a-189e1a8f}{traditional-catholic}{the fall of Jerusalem (587 B.C.)}{Schets, Catholic Encyclopedia}{B.C. 587 (Schets, Catholic Encyclopedia)}, \chronologyannotationanchorclaim{israel.exile.third-captivity}{3}{catholic-traditional-v1}{alternate}{artifact.catholic-encyclopedia.volume-8.new-york-1910.newadvent-08654a-645bba6c}{reported-traditional}{B.C. 588}{Petavius table reported by Sloet; retained transcription checked, printed-page verification still pending}{B.C. 588 (Petavius table reported by Sloet; retained transcription checked, printed-page verification still pending)}. The represented event is after this historical event. These dates belong to that historical event, not to the passage or its composition. Both events could occur in the same year.}%
}
% communion: Communion
\expandafter\def\csname triptychchronologyannotation@communion\endcsname{%
  \chronologyannotationgroup{traditional-attribution}{disputed}{\textbf{Traditional attribution} -- disputed: \chronologyannotationreach{Ps.118.49}{Ps.118}{true}\chronologyannotationreach{Ps.118.50}{Ps.118}{true}\chronologyannotationclaim{israel.monarchy.david-traditional-era}{traditional-attribution}{catholic-traditional-v1}{disputed}{reigned from 1055 to 1015 B.C.}{David (reign in the usual chronology), B.C. 1055--1015}.}%
  \space \chronologyannotationgroup{historical-setting}{research-pending}{\textbf{Historical setting} -- Occasion date unresolved.}%
  \space \chronologyannotationgroup{composition}{preferred}{\textbf{Composition}: \chronologyannotationreach{Ps.118.49}{Ps}{true}\chronologyannotationreach{Ps.118.50}{Ps}{true}\chronologyannotationclaim{critical.psalms.latest-composition-boundary}{composition}{catholic-critical-v1}{preferred}{before the Maccabean period, around 165 B.C.; no individual psalm can be dated securely}{Before c. 165 B.C.}}%
}

\begin{concisedossiertable}
Alleluia & Ps.~107:2 (108) & Davidic title; Israel's praise & \chronodate{alleluia}{\chronologyannotation{alleluia}}\\
\dossierprose{Corbett supplies David's reference era, not this poem's occasion or writing site. USCCB's Psalter-wide boundary dates no individual psalm securely.}
\midrule
Introit verse & Ps.~118:1 (119) & Israel under the law; received Davidic attribution & \chronodate{introit}{\chronologyannotation{introit}}\\
\dossierprose{Ps~118:1: B.C.~1055--1015 is the traditional-attribution regnal reference; Ps~118:1: before c.~165~B.C. is the broad composition boundary. Haydock preserves David instructing Solomon, David persecuted by Saul, and Daniel consoling captives; he favors David. No writing occasion or site is settled.}
\midrule
Communion & Ps.~118:49--50 (119) & A servant's humiliation; same authorship alternatives & \chronodate{communion}{\chronologyannotation{communion}}\\
\dossierprose{Haydock's alternatives apply to the whole poem. Neither the regnal reference nor the broad Psalter boundary dates this humiliation.}
\midrule
Offertory & Ps.~136:1 (137) & Babylonian captives remember Jerusalem & \chronodate{offertory}{\chronologyannotation{offertory}}\\
\dossierprose{Haydock preserves variant headings and captivity, return and prophetic-composition proposals. Writing date/site remain distinct from the represented lament; Jeremiah~25 supplies servitude's duration.}
\midrule
Gradual & Ps.~144:15--16 (145) & Davidic title; Israel praises the Creator & \chronodate{gradual}{\chronologyannotation{gradual}}\\
\dossierprose{No episode or writing site is specified. Corbett dates the attributed king; USCCB supplies the Psalter boundary. Eucharistic reception is later.}
\midrule
Introit antiphon & Dan.~3:29,31,35 & Azarias in Babylon's furnace & Composition: most likely long after the Exile (Gigot).\\
\dossierevent{Daniel~3:23--25 locates the prayer in the furnace. Babylonian servitude is the setting; its ``seventy years'' is a duration, not this episode's date.}
\dossierprose{Gigot reports early attribution but favors an unknown postexilic Jewish author for this portion, distinct from whole-book theories; no precise writing site or life stage is established.}
\midrule
Gospel & John~4:46b--53 & Jesus: Cana; child: Capharnaum, Galilee. Writing: Ephesus traditionally & \chronodate{gospel}{\chronologyannotation{gospel}}\\
\dossierevent{Cana is named in unappointed 4:46a. Maas's encompassing journey yields the derived event range, not a precise healing date or modern clock-time.}
\dossierprose{Fonck and Durand supply composition alternatives for John's later life and a later audience; neither dates the event by the writing.}
\midrule
Epistle & Eph.~5:15--21 & Paul in captivity: Rome favored over Caesarea; Asian recipients & \chronodate{epistle}{\chronologyannotation{epistle}}\\
\dossierprose{Ladeuze and Prat supply traditional alternatives. USCCB distinguishes Paul, a secretary, and a later disciple: style, vocabulary, ecclesiology and Colossians inform debate. The critical interval concerns the disciple only. Impersonal address and omitted place-name witnesses leave destination disputed.}
\end{concisedossiertable}
\zlabel{triptych:concise:chronology:end}
\clearpage
